\documentclass[a4paper]{article}
\usepackage[margin=0.45in,top=0.38in,bottom=0.38in]{geometry}
\usepackage{multicol}
\usepackage{enumitem}
\usepackage{fancyhdr}
\usepackage{xcolor}
\usepackage{tabularx}
\pagestyle{fancy}\fancyhf{}\renewcommand{\headrulewidth}{0pt}\renewcommand{\footrulewidth}{0pt}
\setlength{\columnsep}{0.30in}
\setlength{\parindent}{0pt}\setlength{\parskip}{0pt}
\setlist[enumerate]{leftmargin=1.45em,label=\bfseries\arabic*.,labelsep=0.30em,itemsep=4pt,topsep=0pt,parsep=0pt,partopsep=0pt}
\newcommand{\examheader}{%
\begin{center}\textbf{\large SUMMER EXAMINATION 2026}\\[-1pt]
\textbf{\normalsize Machine Learning}\\[-1pt]
\textbf{BSIT - 3rd Semester}\quad Time: 90 Minutes\quad Max Marks: 80\\[-1pt]
\textcolor{black!65}{\rule{\linewidth}{0.45pt}}\\[-2pt]
\begin{tabularx}{\linewidth}{@{}X X X@{}}Name: \hrulefill & Roll No.: \hrulefill & Date: \hrulefill\end{tabularx}\\[-1pt]
\textcolor{black!65}{\rule{\linewidth}{0.45pt}}\end{center}}
\begin{document}
\examheader
\raggedcolumns
\begin{multicols*}{2}
\fontsize{8.5}{10}\selectfont
\begin{enumerate}
\item Machine learning learns patterns from data to make\newline\hspace*{0.6em}(A) predictions or decisions \quad (B) only web pages \newline\hspace*{0.6em}(C) fixed hand-coded outputs \quad (D) network cables
\item Supervised learning requires\newline\hspace*{0.6em}(A) labeled examples \quad (B) only unlabeled data \newline\hspace*{0.6em}(C) no target variable \quad (D) random rewards
\item Unsupervised learning searches for\newline\hspace*{0.6em}(A) structure without labels \quad (B) known class names \newline\hspace*{0.6em}(C) payment status \quad (D) manual rules
\item Reinforcement learning uses actions, states, and\newline\hspace*{0.6em}(A) rewards \quad (B) SQL joins \newline\hspace*{0.6em}(C) parse trees \quad (D) static labels only
\item The target in a regression problem is usually\newline\hspace*{0.6em}(A) continuous \quad (B) a cluster ID only \newline\hspace*{0.6em}(C) a token sequence \quad (D) a Boolean rule only
\item A feature is an input variable used by a\newline\hspace*{0.6em}(A) model \quad (B) compiler \newline\hspace*{0.6em}(C) database index only \quad (D) test runner
\item A training set is used to\newline\hspace*{0.6em}(A) fit model parameters \quad (B) choose final policy only \newline\hspace*{0.6em}(C) report unbiased test error \quad (D) encrypt records
\item A test set should be held out for\newline\hspace*{0.6em}(A) final generalization evaluation \quad (B) feature creation \newline\hspace*{0.6em}(C) repeated tuning \quad (D) data collection only
\item Missing numeric values may be imputed with the\newline\hspace*{0.6em}(A) median \quad (B) class name only \newline\hspace*{0.6em}(C) file path \quad (D) gradient sign
\item An outlier is an observation that is\newline\hspace*{0.6em}(A) unusually distant from typical values \quad (B) always incorrect \newline\hspace*{0.6em}(C) a missing label \quad (D) a duplicate key only
\item Standardization commonly gives a feature mean 0 and\newline\hspace*{0.6em}(A) standard deviation 1 \quad (B) range 1 \newline\hspace*{0.6em}(C) median 1 \quad (D) variance 0
\item Min-max scaling typically maps values to\newline\hspace*{0.6em}(A) a chosen interval such as 0 to 1 \quad (B) integer labels \newline\hspace*{0.6em}(C) random signs \quad (D) probabilities always
\item One-hot encoding represents a category with\newline\hspace*{0.6em}(A) binary indicator columns \quad (B) a single ordered rank \newline\hspace*{0.6em}(C) a text paragraph \quad (D) a distance matrix only
\item Feature selection removes\newline\hspace*{0.6em}(A) unhelpful input variables \quad (B) all target labels \newline\hspace*{0.6em}(C) test cases \quad (D) model predictions
\item Data leakage occurs when information from the future or target\newline\hspace*{0.6em}(A) enters training improperly \quad (B) is normalized \newline\hspace*{0.6em}(C) is shuffled \quad (D) is missing
\item Linear regression models a target as a\newline\hspace*{0.6em}(A) weighted linear combination \quad (B) tree depth \newline\hspace*{0.6em}(C) cluster assignment \quad (D) probability table only
\item Mean squared error averages\newline\hspace*{0.6em}(A) squared residuals \quad (B) absolute labels \newline\hspace*{0.6em}(C) class probabilities only \quad (D) feature counts
\item Gradient descent moves parameters opposite the\newline\hspace*{0.6em}(A) gradient of loss \quad (B) data mean \newline\hspace*{0.6em}(C) test set \quad (D) class label
\item The learning rate controls\newline\hspace*{0.6em}(A) step size in optimization \quad (B) number of classes \newline\hspace*{0.6em}(C) test proportion only \quad (D) feature count
\item Logistic regression commonly outputs\newline\hspace*{0.6em}(A) a class probability \quad (B) a cluster centroid \newline\hspace*{0.6em}(C) a parse tree \quad (D) a continuous image only
\item The sigmoid function maps real values to\newline\hspace*{0.6em}(A) between 0 and 1 \quad (B) between -1 and 1 only \newline\hspace*{0.6em}(C) all integers \quad (D) unit vectors only
\item KNN predicts using the labels of nearby\newline\hspace*{0.6em}(A) training points \quad (B) decision nodes \newline\hspace*{0.6em}(C) principal components only \quad (D) random samples
\item A larger K in KNN usually makes the decision boundary\newline\hspace*{0.6em}(A) smoother \quad (B) more jagged \newline\hspace*{0.6em}(C) undefined \quad (D) linear always
\item Naive Bayes assumes features are conditionally independent given the\newline\hspace*{0.6em}(A) class \quad (B) distance \newline\hspace*{0.6em}(C) tree depth \quad (D) reward
\item Bayes theorem relates posterior, likelihood, prior, and\newline\hspace*{0.6em}(A) evidence \quad (B) variance only \newline\hspace*{0.6em}(C) gradient \quad (D) residual
\item A decision tree splits data using\newline\hspace*{0.6em}(A) feature tests \quad (B) random rewards \newline\hspace*{0.6em}(C) vector norms only \quad (D) SQL transactions
\item Entropy measures\newline\hspace*{0.6em}(A) impurity or uncertainty \quad (B) distance to origin \newline\hspace*{0.6em}(C) learning rate \quad (D) number of layers
\item Pruning a tree helps reduce\newline\hspace*{0.6em}(A) overfitting \quad (B) missing values \newline\hspace*{0.6em}(C) class count \quad (D) feature scaling
\item An SVM seeks a separating hyperplane with a large\newline\hspace*{0.6em}(A) margin \quad (B) entropy \newline\hspace*{0.6em}(C) learning rate \quad (D) cluster count
\item Support vectors are training points closest to the\newline\hspace*{0.6em}(A) decision boundary \quad (B) centroid \newline\hspace*{0.6em}(C) root node \quad (D) mean label
\item A kernel can enable SVMs to model\newline\hspace*{0.6em}(A) nonlinear boundaries \quad (B) missing values \newline\hspace*{0.6em}(C) test leakage \quad (D) class names
\item Bagging trains models on\newline\hspace*{0.6em}(A) bootstrap samples \quad (B) one fixed sample only \newline\hspace*{0.6em}(C) ordered labels \quad (D) future data
\item Random forests add randomness by selecting subsets of\newline\hspace*{0.6em}(A) features at splits \quad (B) test labels \newline\hspace*{0.6em}(C) reward values \quad (D) output classes only
\item Boosting builds an ensemble\newline\hspace*{0.6em}(A) sequentially focusing on errors \quad (B) in parallel with no interaction \newline\hspace*{0.6em}(C) from one tree \quad (D) only from test data
\item Gradient boosting fits new learners to\newline\hspace*{0.6em}(A) current residual errors \quad (B) feature names \newline\hspace*{0.6em}(C) random labels \quad (D) test predictions only
\item K-means alternates assignment and\newline\hspace*{0.6em}(A) centroid update \quad (B) tree pruning \newline\hspace*{0.6em}(C) gradient clipping \quad (D) label encoding
\item The K in K-means is the number of\newline\hspace*{0.6em}(A) clusters \quad (B) features \newline\hspace*{0.6em}(C) iterations only \quad (D) classes in the test set
\item Hierarchical clustering produces a\newline\hspace*{0.6em}(A) dendrogram \quad (B) confusion matrix \newline\hspace*{0.6em}(C) ROC curve \quad (D) learning curve only
\item DBSCAN identifies dense regions using neighborhood radius and\newline\hspace*{0.6em}(A) minimum points \quad (B) tree depth \newline\hspace*{0.6em}(C) class weights \quad (D) learning rate
\item PCA finds orthogonal directions ordered by\newline\hspace*{0.6em}(A) explained variance \quad (B) class frequency \newline\hspace*{0.6em}(C) reward \quad (D) accuracy only
\end{enumerate}
\end{multicols*}
\newpage
\examheader
\raggedcolumns
\begin{multicols*}{2}
\fontsize{8.5}{10}\selectfont
\begin{enumerate}[resume]
\item The first principal component has the greatest\newline\hspace*{0.6em}(A) variance \quad (B) mean label \newline\hspace*{0.6em}(C) test error \quad (D) number of samples
\item Accuracy can be misleading when classes are\newline\hspace*{0.6em}(A) imbalanced \quad (B) balanced \newline\hspace*{0.6em}(C) continuous \quad (D) standardized
\item Precision equals TP divided by\newline\hspace*{0.6em}(A) TP plus FP \quad (B) TP plus FN \newline\hspace*{0.6em}(C) TN plus FP \quad (D) all negatives
\item Recall equals TP divided by\newline\hspace*{0.6em}(A) TP plus FN \quad (B) TP plus FP \newline\hspace*{0.6em}(C) TN plus FN \quad (D) all predictions
\item For TP=12 and FN=3, recall is\newline\hspace*{0.6em}(A) 0.80 \quad (B) 0.20 \newline\hspace*{0.6em}(C) 0.75 \quad (D) 0.60
\item Specificity is the true negative rate, TN divided by\newline\hspace*{0.6em}(A) TN plus FP \quad (B) TN plus FN \newline\hspace*{0.6em}(C) TP plus FP \quad (D) all positives
\item MAE averages the absolute\newline\hspace*{0.6em}(A) prediction errors \quad (B) squared features \newline\hspace*{0.6em}(C) class labels \quad (D) gradient steps
\item RMSE is the square root of\newline\hspace*{0.6em}(A) MSE \quad (B) MAE \newline\hspace*{0.6em}(C) R squared \quad (D) variance only
\item R squared measures variance explained relative to a\newline\hspace*{0.6em}(A) baseline model \quad (B) confusion matrix \newline\hspace*{0.6em}(C) cluster center \quad (D) learning rate
\item Cross-validation repeatedly splits data into training and\newline\hspace*{0.6em}(A) validation folds \quad (B) test labels only \newline\hspace*{0.6em}(C) features and targets permanently \quad (D) random classes
\item Stratified splitting preserves\newline\hspace*{0.6em}(A) class proportions \quad (B) feature means exactly \newline\hspace*{0.6em}(C) time order always \quad (D) duplicate rows
\item Overfitting is associated with low training error and high\newline\hspace*{0.6em}(A) validation error \quad (B) training accuracy only \newline\hspace*{0.6em}(C) bias only \quad (D) regularization strength
\item Underfitting often indicates a model is\newline\hspace*{0.6em}(A) too simple \quad (B) too complex \newline\hspace*{0.6em}(C) leaking data \quad (D) using too many labels
\item L1 regularization encourages\newline\hspace*{0.6em}(A) sparse coefficients \quad (B) large weights \newline\hspace*{0.6em}(C) more clusters \quad (D) higher variance
\item L2 regularization penalizes the\newline\hspace*{0.6em}(A) sum of squared weights \quad (B) number of classes \newline\hspace*{0.6em}(C) absolute residual only \quad (D) test size
\item A hyperparameter is set\newline\hspace*{0.6em}(A) before or outside parameter fitting \quad (B) only by backpropagation \newline\hspace*{0.6em}(C) by the test label \quad (D) after deployment only
\item Grid search evaluates\newline\hspace*{0.6em}(A) specified hyperparameter combinations \quad (B) all possible datasets \newline\hspace*{0.6em}(C) only one model \quad (D) random gradients
\item A perceptron is a\newline\hspace*{0.6em}(A) single linear threshold unit \quad (B) clustering algorithm \newline\hspace*{0.6em}(C) tree ensemble \quad (D) language model
\item An activation function adds\newline\hspace*{0.6em}(A) nonlinearity \quad (B) database indexing \newline\hspace*{0.6em}(C) data leakage \quad (D) a test split
\item ReLU outputs max of zero and\newline\hspace*{0.6em}(A) its input \quad (B) its square \newline\hspace*{0.6em}(C) its class \quad (D) its probability
\item Backpropagation computes gradients using the\newline\hspace*{0.6em}(A) chain rule \quad (B) Bayes rule only \newline\hspace*{0.6em}(C) Markov rule \quad (D) Pythagorean theorem
\item A CNN convolution filter detects\newline\hspace*{0.6em}(A) local spatial patterns \quad (B) only class priors \newline\hspace*{0.6em}(C) database keys \quad (D) sequence rewards
\item Pooling usually reduces\newline\hspace*{0.6em}(A) spatial resolution \quad (B) number of labels \newline\hspace*{0.6em}(C) input types \quad (D) training epochs only
\item An RNN processes inputs while maintaining a\newline\hspace*{0.6em}(A) hidden state \quad (B) decision stump \newline\hspace*{0.6em}(C) centroid \quad (D) confusion table
\item LSTM gates regulate information in the\newline\hspace*{0.6em}(A) cell state \quad (B) test set \newline\hspace*{0.6em}(C) feature scaler \quad (D) tree root
\item In Q-learning, Q(s,a) estimates expected return for\newline\hspace*{0.6em}(A) an action in a state \quad (B) a feature value \newline\hspace*{0.6em}(C) a class label \quad (D) a cluster
\item An epsilon-greedy policy explores randomly with probability\newline\hspace*{0.6em}(A) epsilon \quad (B) one minus epsilon always \newline\hspace*{0.6em}(C) the learning rate only \quad (D) the discount factor only
\item A discount factor near zero emphasizes\newline\hspace*{0.6em}(A) immediate rewards \quad (B) distant rewards \newline\hspace*{0.6em}(C) feature scaling \quad (D) test accuracy
\item Model drift means the data or relationship changes after\newline\hspace*{0.6em}(A) deployment \quad (B) normalization \newline\hspace*{0.6em}(C) training split \quad (D) label encoding
\item Monitoring should track performance, latency, and\newline\hspace*{0.6em}(A) input or output drift \quad (B) font size \newline\hspace*{0.6em}(C) source comments \quad (D) tree names
\item A reproducible pipeline fixes seeds and records\newline\hspace*{0.6em}(A) data, code, and configuration \quad (B) only accuracy \newline\hspace*{0.6em}(C) only screenshots \quad (D) only labels
\item Fairness evaluation checks whether outcomes differ across\newline\hspace*{0.6em}(A) relevant groups \quad (B) random file names \newline\hspace*{0.6em}(C) feature columns only \quad (D) epochs only
\item Explainability helps stakeholders understand\newline\hspace*{0.6em}(A) model predictions \quad (B) network cables \newline\hspace*{0.6em}(C) database backups \quad (D) file compression
\item A 2 by 2 confusion matrix contains\newline\hspace*{0.6em}(A) four outcome cells \quad (B) two models \newline\hspace*{0.6em}(C) two features \quad (D) four classes necessarily
\item If a classifier gets 90 correct out of 100, accuracy is\newline\hspace*{0.6em}(A) 0.90 \quad (B) 0.10 \newline\hspace*{0.6em}(C) 0.50 \quad (D) 90.0
\item A lower validation loss generally indicates\newline\hspace*{0.6em}(A) better validation fit \quad (B) more leakage necessarily \newline\hspace*{0.6em}(C) higher bias always \quad (D) fewer samples only
\item The residual in regression is actual value minus\newline\hspace*{0.6em}(A) predicted value \quad (B) feature mean \newline\hspace*{0.6em}(C) class prior \quad (D) tree depth
\item The ROC AUC summarizes ranking quality across\newline\hspace*{0.6em}(A) classification thresholds \quad (B) learning rates only \newline\hspace*{0.6em}(C) cluster counts \quad (D) feature encodings
\item Early stopping can reduce overfitting by halting when validation performance\newline\hspace*{0.6em}(A) stops improving \quad (B) reaches zero training loss \newline\hspace*{0.6em}(C) has more features \quad (D) uses no labels
\item An epoch is one pass through the\newline\hspace*{0.6em}(A) training data \quad (B) test labels only \newline\hspace*{0.6em}(C) hyperparameter grid \quad (D) confusion matrix
\end{enumerate}
\end{multicols*}
\end{document}
